\statement{B.80}{Refuted}{One exceptional retained tangency does not force relative source surjectivity}
The following proposed relative special-section principle is false:
if a source line has only one exceptional intersection with the retained
boundary, that exception is a retained tangency, and its intersections
with the deleted divisor and with infinity are transverse, then the
image of its deletion kernel contains the ambient deletion kernel.
The failure already occurs for a parabola and one tangent line, in a
constant-degree affine pencil with a far section. This failure persists for the actual
discrete images and the actual compact images in full profinite
completions. This refutes only the stated general principle; it does
not decide the fifth-root marginal source of B.79.

\begin{proof}
\emph{The actual complements and their source.}\par
In the affine plane with coordinates \((x,y)\), retain the parabola
\(D=\{y=x^2\}\), and delete additionally the line
\(L=\{y=2x-1\}\). The line at infinity is retained throughout.
Put
\[
\begin{gathered}
 Y_0=\mathbb C^2\setminus D,
 \qquad Y=\mathbb C^2\setminus(D\cup L),\\
 S=\{y=0\}\cap Y\simeq
           \mathbb C\setminus\{0,\tfrac12\},
 \qquad
 S_0=\{y=0\}\cap Y_0\simeq\mathbb C\setminus\{0\}.
\end{gathered}
 \tag{B.80.1}
\]
Use the common basepoint \((-1,0)\in S\). Define the actual groups and images by
\[
\begin{gathered}
 G=\pi_1(Y),\qquad
 R=\ker\bigl(G\longrightarrow\pi_1(Y_0)\bigr),\\
 N=\ker\bigl(\pi_1(S)\longrightarrow\pi_1(S_0)\bigr),
 \qquad I=\operatorname{im}(N\longrightarrow G).
\end{gathered}
 \tag{B.80.2}
\]
The source meets the retained parabola only at \((0,0)\), with
intersection multiplicity two. It meets the deleted line transversely
at \((\tfrac12,0)\). In homogeneous coordinates \([X:Y:Z]\), the
source meets infinity at \([1:0:0]\); the parabola \(YZ=X^2\) meets
infinity at \([0:1:0]\), and the deleted line \(Y=2X-Z\) meets it
at \([1:2:0]\). Thus the source has only the one exceptional retained
intersection, while its infinity point is transverse to the retained
infinity line alone.

The pencil \(y=t\) has generic finite punctures
\(\sqrt t,-\sqrt t,(t+1)/2\), with constant total degree three.
There is no deleted entire fiber and no finite escape value. For example,
\((x,y)=(|t|+3,t)\) is a continuous far section, since its first
coordinate is larger in modulus than all three punctures. The other
exceptional value is \(t=1\), where the deleted line is tangent to
the retained parabola. This tangency is away from the source and is
part of the specified example; no assertion about globally transverse
retained and deleted divisors is being made.

\emph{A connected actual eight-sheeted cover.}\par
Set
\[
          u=x-1,\qquad v=y-2x+1.
\]
Then
\[
                Y=\{(u,v):v\ne0,\ v\ne u^2\}.
 \tag{B.80.3}
\]
On the connected line-arrangement complement
\[
 W=\{(\alpha,\beta)\in\mathbb C^2:
              \alpha\beta(\alpha-\beta)(\alpha+\beta)\ne0\}
\]
consider the polynomial map
\[
       u=\frac{\alpha^2+\beta^2}{2},\qquad
       v=\frac{(\alpha^2-\beta^2)^2}{4}.
 \tag{B.80.4}
\]
It maps into \(Y\), because
\(u^2-v=\alpha^2\beta^2\). The two values
\(\alpha^2,\beta^2\) are \(u+\sqrt v,u-\sqrt v\), in either
order. They are distinct and nonzero on \(Y\), and each has two
square roots. Hence every fiber has exactly eight points. Both
\(\alpha\) and \(\beta\) satisfy the monic equation
\[
              Z^4-2uZ^2+u^2-v=0,
\]
so the map is finite. Its Jacobian is
\[
 \det\frac{\partial(u,v)}{\partial(\alpha,\beta)}
       =-2\alpha\beta(\alpha^2-\beta^2),
 \tag{B.80.5}
\]
which never vanishes on \(W\). Thus (B.80.4) is a connected
unramified covering of degree eight. Connectedness follows also
directly by perturbing polygonal paths away from the four complex
lines removed from \(\mathbb C^2\).

The signed coordinate permutations preserve (B.80.4) and act freely
and transitively on each fiber. They form the dihedral group of
order eight, which we denote by
\[
 D_8=\langle d,\sigma\rangle,
 \qquad
 d=\begin{pmatrix}-1&0\\0&1\end{pmatrix},
 \quad
 \sigma=\begin{pmatrix}0&1\\1&0\end{pmatrix}.
 \tag{B.80.6}
\]
Choosing a point over the common source basepoint therefore gives
an actual surjective monodromy homomorphism
\(\rho:G\twoheadrightarrow D_8\). This quotient comes from the
displayed connected cover of the actual surface.

\emph{The exact source components and their monodromy.}\par
The source equation in (B.80.3) is \(v+2u+1=0\). Its pullback is
\[
       \bigl((\alpha+\beta)^2+2\bigr)
       \bigl((\alpha-\beta)^2+2\bigr)=0.
 \tag{B.80.7}
\]
Choose \(c=i\sqrt2\), so \(c^2=-2\). The four components
are the lines
\(\alpha+\beta=\pm c\) and \(\alpha-\beta=\pm c\), with
their intersections with the removed arrangement deleted. Distinct
components meet only at removed points. On the chosen component
\(\alpha-\beta=c\), put \(\beta=t\), \(\alpha=t+c\). Its
actual map to the source is
\[
\begin{gathered}
 x=t(t+c),\\
 \mathbb C\setminus\{0,-c,-c/2\}
       \longrightarrow\mathbb C\setminus\{0,\tfrac12\}.
\end{gathered}
 \tag{B.80.8}
\]
This is a connected unramified double cover of the punctured source.
The two points over \(x=0\) are \(t=0,-c\), both unramified
before removing them. The point over \(x=\tfrac12\) is
\(t=-c/2\), of multiplicity two. The nonidentity deck map of this
component is
\[
             t\longmapsto-c-t,
 \qquad (\alpha,\beta)\longmapsto(-\beta,-\alpha)
                         =-\sigma(\alpha,\beta).
 \tag{B.80.9}
\]
Equivalently, the stabilizer of this source component in \(D_8\)
is exactly \(\{1,-\sigma\}\). To make the cover marking literal,
choose its point over \((-1,0)\) to be
\[
 \alpha_0=\frac{i(\sqrt6+\sqrt2)}2,\qquad
 \beta_0=\frac{i(\sqrt6-\sqrt2)}2.
\]
It has \(\alpha_0-\beta_0=c\) and
\(\beta_0(\beta_0+c)=-1\).

A based positive source meridian around the retained puncture
\(x=0\) has trivial image under \(\rho\), while a based meridian
\(\delta\) around the deleted puncture \(x=\tfrac12\) has image
\(-\sigma\). These two meridians generate \(\pi_1(S)\).
Since \(N\) is their source normal subgroup generated by
\(\delta\), it follows that
\[
        \rho\bigl(\operatorname{im}\pi_1(S)\bigr)
            =\rho(I)=\langle-\sigma\rangle\simeq C_2.
 \tag{B.80.10}
\]
In particular this calculation retains the actual source map and
its deletion kernel, not just a proposed finite peripheral tuple.

\emph{The ambient deletion image is strictly larger.}\par
The part of \(L\) lying in \(Y_0\) is connected: it is the
affine tangent line with its tangency point \((1,1)\) removed.
Transverse perturbation of loops and null-homotopy disks shows
that \(G\to\pi_1(Y_0)\) is onto and that its kernel is the
normal closure of a positive meridian of this divisor. The actual
source meridian \(\delta\) is such a meridian, since the source
meets \(L\) transversely. Surjectivity of \(\rho\) therefore gives
\[
\begin{aligned}
 \rho(R)
   &=\langle\!\langle-\sigma\rangle\!\rangle_{D_8}\\
   &=\{1,-1,\sigma,-\sigma\}\simeq C_2\times C_2,
 \qquad d(-\sigma)d^{-1}=\sigma.
\end{aligned}
 \tag{B.80.11}
\]
Equations (B.80.10)--(B.80.11) prove \(I\ne R\). For example,
a conjugate of \(\delta\) by any actual loop whose \(\rho\)-image
is \(d\) belongs to \(R\) and has image \(\sigma\notin\rho(I)\).

For full profinite completions, define the source image using
\(\ker(\widehat{\pi_1(S)}\to\widehat{\pi_1(S_0)})\), and
define the ambient kernel using
\(\ker(\widehat G\to\widehat{\pi_1(Y_0)})\). The filling
presentations identify these kernels with the respective closed
normal closures of the same meridian. The actual compact source
image is the closure of \(I\), and the ambient closed kernel is
the closure of \(R\). The finite quotient \(\rho\) extends to
\(\widehat G\) and still has the two different images in
(B.80.10)--(B.80.11). Thus it separates these actual compact
images as well; no completion of an unproved infinite-index
subgroup embedding is used.

The independent exact symbolic and eight-element check is
\begin{quote}\small
\path{code/root_checks/check_one_exception_tangent_counterexample.py}.
\end{quote}
Its genuine stdout and
\path{records/one_exception_tangent_counterexample_certificate.json}
verify the coordinate identities, Jacobian, four source components,
component stabilizer, and normal closure. The geometric cover and
the fundamental-group argument above establish the actual-source
claims independently of a bounded search.

The specific fifth-root marginal has the additional internal
normal-generation property inherited from LM4 F.21. The present
\(C_2<V_4\) example does not have that property. Consequently this
counterexample removes the proposed general geometric shortcut
without refuting the marginal equality in B.79 or the original
centralizer problem.
\end{proof}
